\usepackage{tikz}
\usetikzlibrary{positioning, shapes, arrows, shadows, fit, calc, decorations.pathreplacing}

% Semantic Colors - One meaning per color
\definecolor{irSparse}{HTML}{3498DB}   % Sparse/BM25: Blue
\definecolor{irDense}{HTML}{E67E22}    % Dense/Embeddings: Orange
\definecolor{irRerank}{HTML}{C0392B}   % Rerank/Cross-encoder: Red
\definecolor{irHybrid}{HTML}{2ECC71}   % Hybrid/Fusion: Green
\definecolor{irANN}{HTML}{9B59B6}      % ANN/Infra: Purple
\definecolor{irInfra}{HTML}{58595B}    % General Infra: Dark Gray
\definecolor{irMuted}{HTML}{D9D9D9}    % Faint Gray

% Base Styles - Consistent Diagram Grammar
\tikzset{
    irBoxBase/.style={
        rectangle, 
        draw, 
        rounded corners=2pt, 
        minimum height=0.8cm, 
        minimum width=2cm,
        font=\small\sffamily,
        align=center,
        drop shadow={opacity=0.15},
        inner sep=5pt
    },
    irDatabase/.style={
        cylinder, 
        draw, 
        shape border rotate=90, 
        aspect=0.25, 
        minimum height=1.2cm, 
        minimum width=1cm,
        fill=irMuted!20,
        draw=irInfra,
        font=\tiny\sffamily,
        align=center
    },
    % Semantic Variants
    sparse/.style={irBoxBase, fill=irSparse!10, draw=irSparse},
    dense/.style={irBoxBase, fill=irDense!10, draw=irDense},
    rerank/.style={irBoxBase, fill=irRerank!10, draw=irRerank},
    hybrid/.style={irBoxBase, fill=irHybrid!10, draw=irHybrid},
    ann/.style={irBoxBase, fill=irANN!10, draw=irANN},
    infra/.style={irBoxBase, fill=irInfra!10, draw=irInfra},
    % Legacy support / Utility
    muted/.style={irBoxBase, fill=irMuted!20, draw=irMuted!80!black, text=irMuted!50!black},
    ghost/.style={irBoxBase, fill=white, draw=irMuted!50, text=irMuted, dashed, drop shadow={opacity=0}},
    emph/.style={draw=irRerank, thick, drop shadow={shadow xshift=0mm, shadow yshift=0mm, fill=irRerank, opacity=0.3}},
    % Arrows
    irArrowBase/.style={->, thick, >=latex, draw=irInfra},
    irLabel/.style={
        font=\tiny\itshape, 
        text=irInfra, 
        align=center, 
        text width=2.5cm, 
        inner sep=3pt
    }
}


% =========================================================
% Stage Badge System
% =========================================================
% \irStage{StageName}
\newcommand{\irStage}[1]{%
    \begin{tikzpicture}[remember picture, overlay]
        \node[
            anchor=north east,
            fill=irInfra!10,
            text=irInfra,
            font=\tiny\bfseries,
            inner sep=3pt,
            rounded corners=1pt,
            yshift=-0.5cm,
            xshift=-0.5cm
        ] at (current page.north east) {#1};
    \end{tikzpicture}%
}

% =========================================================
% Math Highlighting
% =========================================================
% \mathhl[color]{content}
\newcommand{\mathhl}[2][irDense!30]{%
    \colorbox{#1}{$\displaystyle #2$}%
}

% =========================================================
% Math Vectors
% =========================================================
\newcommand{\vs}{\mathbf{s}}

% =========================================================
% Core Box Primitives
% =========================================================
% \IRBox[width]{id}{at}{title}{subtitle}{variant}
% Default width chosen for lecture slides; override when needed.
% \IRBox[width]{id}{at}{title}{subtitle}{variant}[extra options]
\DeclareDocumentCommand{\IRBox}{O{3.5cm} m m m m m O{}}{%
    \node[#6, text width=#1, #7] (#2) at #3 {%
        \textbf{#4}%
        \if\relax\detokenize{#5}\relax
        \else\\[-1mm]\scriptsize #5%
        \fi
    };%
}

% \IRWideBox{id}{at}{width}{title}{subtitle}{variant}
\newcommand{\IRWideBox}[6]{%
    \node[#6, text width=#3] (#1) at #2 {%
        \textbf{#4}%
        \if\relax\detokenize{#5}\relax
        \else\\[-1mm]\scriptsize #5%
        \fi
    };%
}

% =========================================================
% Arrow Primitives
% =========================================================
% Safe default: left-to-right flow
% \IRArrow{from}{to}{label}
\newcommand{\IRArrow}[3]{%
    \draw[irArrowBase] (#1.east) -- node[irArrowLabelAbove] {#3} (#2.west);%
}

% Explicit left-to-right
% \IRArrowLR{from}{to}{label}
\newcommand{\IRArrowLR}[3]{%
    \draw[irArrowBase] (#1.east) -- node[irArrowLabelAbove] {#3} (#2.west);%
}

% Explicit top-to-bottom
% \IRArrowTB{from}{to}{label}
\newcommand{\IRArrowTB}[3]{%
    \draw[irArrowBase] (#1.south) -- node[irArrowLabelRight] {#3} (#2.north);%
}

% Fully explicit anchor-to-anchor path
% \IRArrowPath{from anchor}{to anchor}{label}
% Example: \IRArrowPath{a.east}{b.north}{score}
\newcommand{\IRArrowPath}[3]{%
    \draw[irArrowBase] (#1) -- node[irArrowLabelAbove] {#3} (#2);%
}

% Orthogonal elbow path
% \IRArrowElbow{from anchor}{to anchor}{label}
% Example: \IRArrowElbow{a.east}{b.north}{feedback}
\newcommand{\IRArrowElbow}[3]{%
    \draw[irArrowBase] (#1) -| node[irArrowLabelBox] {#3} (#2);%
}

% Slanted arrow for special cases only
% \IRArrowSlanted{from anchor}{to anchor}{label}{pos}{width}
\newcommand{\IRArrowSlanted}[5]{%
    \draw[irArrowBase] (#1) -- node[above, sloped, irLabel, pos=#4, text width=#5, fill=white, inner sep=1pt] {#3} (#2);%
}

% =========================================================
% Accent Elements
% =========================================================
% \IRBadge{at}{text}
\newcommand{\IRBadge}[2]{%
    \node[
        draw=irAccent,
        fill=irAccent!20,
        rounded corners=1pt,
        font=\tiny\bfseries,
        inner sep=2pt
    ] at #1 {#2};%
}

% =========================================================
% Grouping
% =========================================================
% \IRGroup{fit nodes}{label}{variant}
% Example: \IRGroup{(a)(b)(c)}{First-stage retrieval}{muted}
\newcommand{\IRGroup}[3]{%
    \node[
        irGroupBase,
        #3,
        fit=#1
    ] (group) {};
    \node[
        anchor=north west,
        font=\tiny\bfseries,
        #3
    ] at ([xshift=2pt,yshift=-2pt] group.north west) {#2};%
}

% =========================================================
% Domain-Specific Shapes
% =========================================================
% \IRDocStack{id}{at}{label}
\newcommand{\IRDocStack}[3]{%
    \node[irDocStack] (#1) at #2 {#3};%
}

% \IRDatabase{id}{at}{label}
\newcommand{\IRDatabase}[3]{%
    \node[irDatabase] (#1) at #2 {#3};%
}

% =========================================================
% Layout Helpers
% =========================================================
% \IRPipelineTier{id}{at}{text}{label}{variant}
\newcommand{\IRPipelineTier}[5]{%
    \IRBox{#1}{#2}{#3}{}{#5}[]
    \node[anchor=west, irLabel, xshift=0.5cm] at (#1.east) {#4};%
}

% \IRTelescopeTier{id}{at}{width}{text}{variant}
\newcommand{\IRTelescopeTier}[5]{%
    \node[
        irTelescopeBase,
        #5,
        minimum width=#3
    ] (#1) at #2 {#4};%
}

% \IRStep{id}{at}{title}{subtitle}{variant}
% Creates a numbered step badge left of a box.
\newcommand{\IRStep}[5]{%
    \node[
        circle,
        draw=irAccent,
        fill=irAccent!20,
        font=\tiny\bfseries,
        inner sep=2pt
    ] (#1_num) at #2 {#1};
    \node[
        #5,
        text width=3.5cm,
        right=1.5cm of #1_num
    ] (#1) {%
        \textbf{#3}%
        \if\relax\detokenize{#4}\relax
        \else\\[-1mm]\scriptsize #4%
        \fi
    };
    \draw[irArrowBase, thin] (#1_num.east) -- (#1.west);%
}
